\section{Introduction}
\label{sec:intro}

The migration to post-quantum cryptography (PQC) is well under way.
Two years after NIST published ML-KEM~\cite{FIPS203}, ML-DSA~\cite{FIPS204}
and SLH-DSA~\cite{FIPS205}, post-quantum key exchange is already widely
deployed, whereas post-quantum authentication lags behind even though its
deadlines are approaching.
NIST's draft transition plan deprecates quantum-vulnerable signatures such as
RSA and ECDSA at 112-bit security after 2030 and disallows them all after
2035~\cite{NISTIR8547}, the EU roadmap recommends migrating high-risk use
cases by the end of 2030~\cite{EUPQC25}, and for U.S.\ national security
systems, the NSA expects software and firmware signing to use post-quantum
algorithms exclusively by 2030~\cite{CNSA2}.
Embedded devices are a particularly hard case: devices deployed today must
verify firmware and authenticate themselves well beyond these dates, under
tight bandwidth, memory and energy budgets.
For them, FN-DSA (FIPS~206), derived from the lattice-based signature scheme
Falcon~\cite{Falcon}, is the most attractive of NIST's stateless post-quantum
signatures: its public key and signature take 897 and 666~bytes for
FN-DSA-512, compared with 1312 and 2420~bytes for ML-DSA-44~\cite{FIPS204},
and its verification needs little more than hashing and a single polynomial
multiplication modulo $q=12289$.
At the time of writing, the FIPS~206 draft has not yet been published, so we
follow the reference implementation~\cite{cfndsa}, which tracks the expected
draft.
\TODO{Re-check all dates, figures and the FIPS~206 status before
submission.}
% [역할] 연구 분야의 큰 흐름: 2026년 PQC 전환 현황(키 교환은 대규모 배포,
%        서명은 지연)과 다가오는 시한 → 임베디드가 가장 어려움 → FN-DSA가 적합
% [번역]
% 양자내성암호(PQC)로의 전환은 이미 한창 진행 중이다.
% NIST가 ML-KEM, ML-DSA, SLH-DSA를 발표한 지 2년이 지난 지금, 양자내성
% 키 교환은 이미 널리 배포된 반면, 양자내성 인증(서명)은 시한이 다가오고
% 있는데도 뒤처져 있다.
% NIST의 전환 계획 초안은 RSA, ECDSA 같은 양자 취약 서명의 112비트
% 보안 수준을 2030년 이후 폐기 예정(deprecated)으로, 2035년 이후에는
% 모두 사용 금지(disallowed)로 정하고 있다. EU 로드맵은 고위험 사용처를
% 2030년 말까지 전환하도록 권고하며, 미국 국가안보시스템(NSS)에 대해서는
% NSA(CNSA 2.0)가 소프트웨어·펌웨어 서명이 2030년까지 양자내성 알고리즘만
% 쓰도록 예정해 두었다.
% 임베디드 기기는 특히 어려운 경우다. 지금 배포되는 기기는 이 시한을
% 훨씬 넘어서까지 펌웨어를 검증하고 자신을 인증해야 하며, 대역폭·메모리·
% 에너지 제약도 빠듯하다.
% 이런 기기에는 격자 기반 서명 Falcon에서 나온 FN-DSA(FIPS 206)가 NIST의
% 무상태(stateless) 양자내성 서명 중 가장 매력적이다. FN-DSA-512의
% 공개키와 서명은 897바이트와 666바이트로 ML-DSA-44의 1312바이트와
% 2420바이트보다 작고, 검증에는 해싱과 q=12289 모듈러 다항식 곱셈 한 번
% 정도만 필요하다.
% 이 글을 쓰는 시점에 FIPS 206 초안은 아직 공개되지 않았으므로, 예상
% 초안을 따라가는 레퍼런스 구현을 기준으로 한다.
% (TODO: 제출 전에 모든 날짜·수치와 FIPS 206 상태 다시 확인)

This compactness comes at the price of a delicate implementation.
Signing relies on fast Fourier sampling (ffSampling)~\cite{DP16,Falcon}, a
randomized nearest-plane algorithm that operates on polynomials in the complex
Fourier domain and repeatedly calls a discrete Gaussian sampler whose center
and standard deviation depend on the secret key.
These computations are specified in IEEE-754 binary64 arithmetic, and their
numerical behavior matters for security, not only for performance:
Lin et al.~\cite{LTYZ25} showed that even tiny deviations between
mathematically equivalent floating-point implementations of the sampler can,
under certain conditions, have security consequences.
A conservative way to rule out such deviations is to require that an
optimized implementation reproduces the known-answer tests (KATs) of the
reference implementation bit for bit.
Moreover, all floating-point code must run in constant time, and side-channel
analysis of Falcon's floating-point arithmetic and Gaussian sampling has
repeatedly led to key recovery~\cite{KA21,GMRR22,ZLYW23}.
Key generation is demanding as well: it solves the NTRU equation with
big-integer arithmetic in a residue number system (RNS) and with Babai
reductions whose approximations are computed in the Fourier domain, using
fixed-point arithmetic in the most efficient
implementations~\cite{Pornin23kg,Pornin25kg}.
% [역할] 문제 정의: FN-DSA 구현이 까다로운 이유 (binary64 의존, 정확성이
%        보안에 미치는 영향과 KAT bit 단위 재현, 상수시간, 키 생성의 복잡도)
% [번역]
% 이 간결함의 대가로 구현은 까다롭다.
% 서명은 fast Fourier sampling(ffSampling)에 의존한다. 이는 복소
% Fourier 영역의 다항식 위에서 동작하는 무작위 nearest-plane
% 알고리즘으로, 비밀키에 따라 중심과 표준편차가 달라지는 이산
% 가우시안 샘플러를 반복해서 호출한다.
% 이 계산들은 IEEE-754 binary64 연산으로 규정되어 있으며, 그 수치적
% 동작은 성능뿐 아니라 보안에도 중요하다. Lin 등은 수학적으로 동등한
% 샘플러의 부동소수점 구현들 사이의 아주 작은 차이도 특정 조건에서는
% 보안 문제로 이어질 수 있음을 보였다.
% 이런 차이를 보수적으로 배제하는 방법은, 최적화 구현이 레퍼런스 구현의
% KAT(known-answer test)를 bit 단위로 그대로 재현하도록 요구하는 것이다.
% 또한 모든 부동소수점 코드는 상수시간으로 동작해야 하며, Falcon의
% 부동소수점 연산과 가우시안 샘플링에 대한 부채널 분석은 여러 차례
% 키 복구로 이어졌다.
% 키 생성도 까다롭다. RNS(잉여수 체계) 위의 큰 정수 연산과 Babai
% reduction으로 NTRU 방정식을 푸는데, Babai reduction의 근사값은
% Fourier 영역에서 계산하며, 가장 효율적인 구현들은 이를 고정소수점
% 연산으로 수행한다.

On microcontrollers, these requirements lead to a difficult trade-off.
The Arm Cortex-M4, the de facto reference platform for embedded PQC, has at
most a single-precision floating-point unit (FPU), so the fastest
implementations of Falcon emulate binary64 arithmetic with integer
instructions~\cite{Pornin19,Pornin25m4}.
Despite years of optimization, signing with FN-DSA-512 still takes more than
ten million cycles on this platform~\cite{Pornin25m4,Pornin26ram}.
Several directions have been explored to reduce this cost.
Howe and Westerbaan benchmarked Falcon on the Cortex-M7, which has a
double-precision FPU~\cite{HW23}.
TWFalcon trades speed for precision with triple-word single-precision
arithmetic on the Cortex-M4F~\cite{TWFalcon}.
Fixed-point arithmetic has been proposed for key
generation~\cite{Pornin23kg,Pornin25kg} and, more recently, for
signing~\cite{BFLP26}.
Finally, Pornin recently removed more than half of the floating-point
operations of signing, partly by replacing them with modular integer
computations, primarily to reduce RAM usage~\cite{Pornin26ram}.
In short, existing microcontroller implementations either emulate binary64,
replace it with alternative arithmetic, or target cores without a vector
unit.
% [역할] 공통 한계: MCU 선행연구의 접근 방식과, 이들의 공통적인 한계 요약
% [번역]
% 마이크로컨트롤러에서는 이런 요구사항 때문에 어려운 절충이 생긴다.
% 임베디드 PQC의 사실상 기준 플랫폼인 Arm Cortex-M4는 기껏해야
% 단정밀도 FPU만 갖추고 있어서, 가장 빠른 Falcon 구현들은 binary64
% 연산을 정수 명령으로 에뮬레이션한다.
% 수년간의 최적화에도 이 플랫폼에서 FN-DSA-512 서명은 여전히 천만
% 사이클 이상 걸린다.
% 이 비용을 줄이기 위해 여러 방향이 시도되었다.
% Howe와 Westerbaan은 배정밀도 FPU를 갖춘 Cortex-M7에서 Falcon을
% 벤치마크했다.
% TWFalcon은 Cortex-M4F에서 삼중 워드(triple-word) 단정밀도 연산으로
% 속도를 내주고 정밀도를 얻었다.
% 고정소수점 연산은 키 생성에, 최근에는 서명에도 제안되었다.
% 마지막으로 Pornin은 최근 주로 RAM 사용량을 줄이기 위해 서명의
% 부동소수점 연산을 절반 넘게 없앴고, 그중 일부는 모듈러 정수 연산으로
% 대체했다.
% 요컨대 기존 MCU 구현은 binary64를 에뮬레이션하거나, 다른 산술로
% 대체하거나, 벡터 유닛이 없는 코어를 대상으로 한다.

The Armv8.1-M architecture changes this landscape.
Its first implementation, the Cortex-M55, and its successor, the Cortex-M85,
combine an FPU that optionally supports double precision with the M-Profile
Vector Extension (MVE), also known as Helium~\cite{ArmV8MARM}.
Both cores ship in commercial microcontrollers, e.g., the Cortex-M55 in the
STM32N6 series~\cite{STM32N6DS}.
This suggests a simple recipe for FN-DSA: run the floating-point parts on the
double-precision FPU and vectorize the integer parts with MVE.
However, the two features interact in ways that make this recipe far from
obvious.
First, MVE offers 128-bit vectors of 8-, 16- and 32-bit integers and of
half- and single-precision floats, but no double-precision lanes; hence, none
of Falcon's binary64 arithmetic can be vectorized directly.
Second, the eight vector registers \texttt{Q0}--\texttt{Q7} alias the
floating-point register file, so scalar FP64 code and MVE code compete for
the same sixteen double-precision registers \texttt{D0}--\texttt{D15}.
Third, the Cortex-M55 executes each 128-bit vector instruction in several
beats and overlaps instructions that use different execution resources, so
its throughput depends strongly on instruction mix and
scheduling~\cite{BMKYV22,SLOTHY}.
As a result, each numerical kernel of FN-DSA admits at least three competing
implementation strategies: scalar binary64 on the FPU, fixed-point or
emulated arithmetic vectorized with integer MVE, and multi-word
single-precision arithmetic vectorized with floating-point MVE.
Since these strategies stress different parts of the microarchitecture, the
right choice has to be made kernel by kernel.
% [역할] 문제의 중요성: M55가 주는 기회와, 단순한 공식이 통하지 않는
%        세 가지 이유 → 커널마다 산술을 골라야 한다는 결론
% [번역]
% Armv8.1-M 아키텍처는 이 상황을 바꾼다.
% 그 첫 구현인 Cortex-M55와 후속 Cortex-M85는 배정밀도를 선택적으로
% 지원하는 FPU와 M-Profile Vector Extension(MVE, Helium)을 함께
% 갖추고 있다.
% 두 코어 모두 상용 MCU로 출시되어 있으며, 예를 들어 STM32N6 시리즈가
% Cortex-M55를 탑재한다.
% 여기서 FN-DSA에 대한 단순한 공식을 떠올릴 수 있다. 부동소수점 부분은
% 배정밀도 FPU에서 실행하고, 정수 부분은 MVE로 벡터화하는 것이다.
% 그러나 두 기능은 서로 얽혀 있어서 이 공식은 결코 자명하지 않다.
% 첫째, MVE는 8·16·32비트 정수와 반정밀도·단정밀도 부동소수점의
% 128비트 벡터를 제공하지만 배정밀도 lane은 없다. 따라서 Falcon의
% binary64 연산은 어느 것도 직접 벡터화할 수 없다.
% 둘째, 8개의 벡터 레지스터 Q0–Q7이 부동소수점 레지스터 파일과
% 겹치므로, 스칼라 FP64 코드와 MVE 코드는 같은 16개의 배정밀도
% 레지스터 D0–D15를 두고 경쟁한다.
% 셋째, Cortex-M55는 128비트 벡터 명령 하나를 여러 beat에 걸쳐
% 실행하고, 서로 다른 실행 자원을 쓰는 명령들을 겹쳐 실행한다. 그래서
% 처리량이 명령 조합과 스케줄링에 크게 좌우된다.
% 결과적으로 FN-DSA의 각 수치 커널에는 적어도 세 가지 구현 전략이
% 경쟁한다. FPU의 스칼라 binary64, 정수 MVE로 벡터화한 고정소수점
% 또는 에뮬레이션 연산, 부동소수점 MVE로 벡터화한 다중 워드 단정밀도
% 연산이다.
% 이 전략들은 마이크로아키텍처의 서로 다른 부분에 부담을 주므로,
% 올바른 선택은 커널마다 따로 내려야 한다.

Existing work leaves two questions open.
First, it offers no guidance on how to divide FN-DSA between a scalar FP64
unit and a vector unit without double-precision lanes.
Vectorized Falcon implementations for NEON, AVX2, AVX-512 and
RVV~\cite{KSS22,NG23,ZZ26} rely on double-precision SIMD instructions, which
MVE lacks.
The Cortex-M7 results of Howe and Westerbaan~\cite{HW23} use a scalar
double-precision FPU, but on a core without a vector extension.
Conversely, Helium implementations of lattice-based schemes, such as Kyber,
Dilithium and Saber~\cite{BMKYV22,SLOTHY} or HAETAE and
SMAUG-T~\cite{LSS26}, target integer and fixed-point arithmetic, since none
of these schemes requires binary64.
Second, it is unclear how much acceleration remains once the outputs must
stay exact.
The most recent vectorized implementation of Falcon signing~\cite{ZZ26}
obtains part of its speedup by batching several calls of the base sampler
and by using four-way parallel SHAKE, both of which make its signatures
differ from those of the reference implementation.
To the best of our knowledge, no optimized implementation of Falcon or FN-DSA
for the Cortex-M55 has been published.
% [역할] 연구 공백(open problem) 두 가지: ① FP64 SIMD가 없는 환경에서의
%        산술 선택 기준이 없음 ② 출력 정확성을 지킬 때 남는 가속 폭을 모름
% [번역]
% 기존 연구는 두 가지 질문을 열어 둔다.
% 첫째, 배정밀도 lane이 없는 벡터 유닛과 스칼라 FP64 유닛 사이에
% FN-DSA를 어떻게 나눌지에 대한 지침이 없다.
% NEON, AVX2, AVX-512, RVV용 벡터화 Falcon 구현들은 MVE에 없는
% 배정밀도 SIMD 명령에 의존한다.
% Howe와 Westerbaan의 Cortex-M7 결과는 스칼라 배정밀도 FPU를 쓰지만,
% 벡터 확장이 없는 코어에서의 결과다.
% 반대로 Kyber, Dilithium, Saber나 HAETAE, SMAUG-T 같은 격자 기반
% 방식의 Helium 구현들은 정수·고정소수점 연산을 대상으로 한다. 이
% 방식들은 binary64를 필요로 하지 않기 때문이다.
% 둘째, 출력이 정확하게 유지되어야 할 때 가속이 얼마나 남는지 불분명하다.
% Falcon 서명의 가장 최근 벡터화 구현은 기본 샘플러의 여러 호출을
% 묶어서 처리하고 4-way 병렬 SHAKE를 사용해 속도 향상의 일부를 얻는데,
% 두 방법 모두 서명이 레퍼런스 구현의 출력과 달라지게 만든다.
% 우리가 아는 한, Cortex-M55용으로 최적화된 Falcon 또는 FN-DSA 구현은
% 아직 발표되지 않았다.

\paragraph{Our approach.}
We present the first optimized implementation of FN-DSA for the Arm
Cortex-M55.
It accelerates every expensive kernel of key generation, signing and
verification for FN-DSA-512 and FN-DSA-1024, and we evaluate it on real
hardware, an STM32N657 microcontroller running at 800\,MHz.
Profiling the reference implementation~\cite{cfndsa} on this board shows
why all three operations need attention: floating-point-heavy code accounts
for about \tbd{87\%} of the signing time, whereas key generation is
dominated by the integer arithmetic of the NTRU solver and verification by
the NTT and SHAKE.
Each of our techniques results from a per-kernel selection: we implement
candidate strategies for a kernel, measure them on the device, discard those
that break the KATs or introduce secret-dependent timing, and adopt the
fastest remaining one.
Applied to FN-DSA, this rule gives a different answer for each class of
kernels.
For the integer kernels, MVE is the natural choice, as expected.
For the Fourier-domain arithmetic of key generation, the FP64 unit loses
against fixed-point arithmetic vectorized with integer MVE.
For signing, the FP64 unit wins and replaces binary64 emulation without
changing a single bit of the output.
\TODO{Confirm after implementation.}
Finally, for the base sampler and SHAKE, exactness rules out the usual way
of using SIMD across independent instances, so we parallelize within a
single instance instead.
% [역할] 제안 기법의 개요: 무엇을 했는지, 기법을 고르는 규칙, 그리고
%        본론 이야기 흐름(5~8절) 미리보기
% [번역]
% 우리는 Arm Cortex-M55를 위한 첫 번째 최적화 FN-DSA 구현을 제시한다.
% 이 구현은 FN-DSA-512와 FN-DSA-1024의 키 생성·서명·검증에서 비용이
% 큰 모든 커널을 가속하며, 800 MHz로 동작하는 실제 하드웨어인
% STM32N657 MCU에서 평가한다.
% 이 보드에서 레퍼런스 구현을 프로파일링하면 세 연산 모두를 개선해야
% 하는 이유가 드러난다. 서명 시간의 약 87%를 부동소수점 연산이 많은
% 코드가 차지하는 반면, 키 생성은 NTRU solver의 정수 연산이, 검증은
% NTT와 SHAKE가 지배한다.
% 우리의 각 기법은 커널별 선택의 결과다. 한 커널의 후보 전략들을
% 구현해 보드에서 측정하고, KAT를 깨거나 비밀값에 따라 실행 시간이
% 달라지는 후보를 버린 뒤, 남은 후보 중 가장 빠른 것을 채택한다.
% 이 규칙을 FN-DSA에 적용하면 커널 부류마다 다른 답이 나온다.
% 정수 커널에서는 예상대로 MVE가 자연스러운 선택이다.
% 키 생성의 Fourier 영역 연산에서는 FP64 유닛이 정수 MVE로 벡터화한
% 고정소수점 연산에 진다.
% 서명에서는 FP64 유닛이 이기며, 출력을 단 한 bit도 바꾸지 않고
% binary64 에뮬레이션을 대체한다. (TODO: 구현 후 확인)
% 마지막으로 기본 샘플러와 SHAKE에서는 정확성 조건 때문에 독립된
% 인스턴스들에 걸쳐 SIMD를 쓰는 일반적인 방법을 쓸 수 없으므로, 대신
% 하나의 인스턴스 안에서 병렬화한다.

\paragraph{Contributions.}
Our contributions are as follows.
% [역할] 논문의 주요 기여: 본론 절(5~8절)에 대응하는 개선 기법과 연구의 부산물
% [번역]
% 우리의 기여는 다음과 같다.
\begin{itemize}
  % Evidence: fn-dsa_m55/ntt_opt*/, fn-dsa_m55/ntt_k4c_compare/result.md
  \item \textbf{Vectorized integer kernels (Section~\ref{sec:ntt}).}
    We vectorize the NTT modulo $q=12289$, which all three operations use,
    and the NTTs over 31-bit primes used by the NTRU solver, choosing
    reduction methods, layer splits and instruction schedules by on-device
    comparison.
    This alone speeds up verification by \tbd{1.16$\times$} and
    \tbd{1.17$\times$} for FN-DSA-512 and FN-DSA-1024, respectively.
    % [번역] 벡터화한 정수 커널 (5절).
    % 세 연산이 모두 쓰는 q=12289 모듈러 NTT와, NTRU solver가 쓰는
    % 31비트 소수 위의 NTT를 벡터화한다. reduction 방식, 레이어 분할,
    % 명령 스케줄은 보드 위 비교로 고른다.
    % 이것만으로 FN-DSA-512와 FN-DSA-1024의 검증이 각각 1.16배,
    % 1.17배 빨라진다.
  % Evidence: fn-dsa_m55/function_compare/{fft_keygen,twfalcon}/,
  %           fn-dsa_m55/sep_27_ntrusolve_fft/result.md
  \item \textbf{Fourier-domain arithmetic of key generation
    (Section~\ref{sec:keygen}).}
    Comparing native binary64, multi-word single-precision MVE and
    integer-MVE fixed-point implementations of the Babai reduction, we find
    that native binary64 makes the transforms \tbd{21.5--30.4\%} slower
    than the reference fixed-point code, whereas integer-MVE fixed-point
    arithmetic is the fastest and yields the same keys as the reference.
    Together with the vectorized NTTs, key generation becomes
    \tbd{1.22$\times$} and \tbd{1.17$\times$} faster.
    % [번역] 키 생성의 Fourier 영역 연산 (6절).
    % Babai reduction을 native binary64, 다중 워드 단정밀도 MVE, 정수 MVE
    % 고정소수점으로 구현해 비교한 결과, native binary64는 변환을
    % 레퍼런스 고정소수점 코드보다 21.5–30.4% 느리게 만드는 반면, 정수
    % MVE 고정소수점 연산은 가장 빠르면서 레퍼런스와 같은 키를 만든다.
    % 벡터화한 NTT와 함께 키 생성이 1.22배, 1.17배 빨라진다.
  \item \textbf{Bit-exact native binary64 for signing
    (Section~\ref{sec:sign}).}
    \TODO{Confirm after implementation.}
    We replace binary64 emulation in the FFT and ffSampling with the scalar
    FP64 unit while reproducing the KATs bit for bit, which multi-word
    single-precision alternatives cannot do.
    % [번역] 서명을 위한 bit 단위로 정확한 native binary64 (7절).
    % (TODO: 구현 후 확인) FFT와 ffSampling의 binary64 에뮬레이션을
    % 스칼라 FP64 유닛으로 대체하면서 KAT를 bit 단위로 그대로 재현한다.
    % 다중 워드 단정밀도 대안들은 이를 할 수 없다.
  \item \textbf{Exactness-preserving vectorization
    (Section~\ref{sec:exact}).}
    \TODO{Confirm after implementation.}
    We vectorize the table comparisons of the base sampler within a single
    sample, a KAT-compatible alternative to batching~\cite{ZZ26} that the
    AVX2 code of the original Falcon submission also
    follows~\cite{FalconImpl}, and we accelerate single-state Keccak, since
    every use of SHAKE in the current reference implementation processes a
    single stream~\cite{cfndsa}.
    % [번역] 정확성을 유지하는 벡터화 (8절).
    % (TODO: 구현 후 확인) 기본 샘플러의 표 비교를 한 샘플 안에서
    % 벡터화한다. 이는 여러 샘플을 묶는 방식[ZZ26]의 KAT 호환 대안이며,
    % 원래 Falcon 제출본의 AVX2 코드도 같은 방식을 쓴다. 또한 현재
    % 레퍼런스 구현에서 SHAKE를 쓰는 모든 곳이 단일 스트림을 처리하므로,
    % 단일 상태 Keccak을 가속한다.
  \item \textbf{By-products.}
    \TODO{Finalize.}
    A timing characterization of the FP64 instructions of the Cortex-M55,
    including subnormal operands; guidelines for choosing between scalar
    FP64 and MVE arithmetic on Armv8.1-M; and our code and measurement
    framework, which are publicly available at
    \TODO{anonymized repository link}.
    % [번역] 부산물. (TODO: 확정)
    % Cortex-M55 FP64 명령의 타이밍 특성 분석(subnormal 피연산자 포함),
    % Armv8.1-M에서 스칼라 FP64와 MVE 연산 중 무엇을 고를지에 대한
    % 지침, 그리고 공개하는 코드와 측정 프레임워크(익명 저장소 링크는
    % 추후 기입).
\end{itemize}
Overall, compared with the reference implementation ported to the same
board, our implementation speeds up key generation, signing and verification
by \tbd{X$\times$}, \tbd{Y$\times$} and \tbd{Z$\times$} for FN-DSA-512, and
by \tbd{X$\times$}, \tbd{Y$\times$} and \tbd{Z$\times$} for FN-DSA-1024.
\TODO{Position against the Cortex-M4~\cite{Pornin25m4,Pornin26ram} and
Cortex-M7~\cite{HW23} results.}
% [역할] 전체 성과 요약과 선행연구 대비 위치
% [번역]
% 종합하면, 같은 보드로 이식한 레퍼런스 구현과 비교해 우리 구현은
% FN-DSA-512의 키 생성·서명·검증을 각각 X배, Y배, Z배,
% FN-DSA-1024에서는 X배, Y배, Z배 빠르게 한다.
% (TODO: Cortex-M4 및 Cortex-M7 결과와 비교한 위치 서술)

\paragraph{Organization.}
Section~\ref{sec:prelim} recalls FN-DSA, its reference implementation and
the relevant features of the Cortex-M55.
Section~\ref{sec:related} discusses related work and derives our design
considerations.
Section~\ref{sec:method} describes our experimental setup, profiles the
reference implementation and states our selection rule.
Sections~\ref{sec:ntt} to~\ref{sec:exact} present our techniques for the
four classes of kernels outlined above.
Section~\ref{sec:eval} evaluates performance, memory usage and constant-time
behavior, Section~\ref{sec:discussion} discusses lessons, limitations and
future work, and Section~\ref{sec:conclusion} concludes.
% [역할] 논문 구성
% [번역]
% Preliminaries 절은 FN-DSA와 그 레퍼런스 구현, Cortex-M55의 관련 특징을
% 다룬다. Related Work 절은 관련 연구를 정리하고 설계 고려사항을
% 도출한다. Design Methodology 절은 실험 환경을 설명하고, 레퍼런스
% 구현을 프로파일링하며, 선택 규칙을 제시한다. 5절부터 8절까지는 위에서
% 요약한 네 부류의 커널을 위한 기법을 제시한다. Evaluation 절은 성능,
% 메모리 사용량, 상수시간 동작을 평가하고, Discussion 절은 교훈, 한계,
% 향후 과제를 논의하며, Conclusion 절에서 결론을 맺는다.
